\documentclass[11pt]{article}
\usepackage[a4paper, top=18mm, bottom=18mm, left=20mm, right=20mm]{geometry}
\usepackage[T2A]{fontenc}
\usepackage[utf8]{inputenc}
\usepackage[ukrainian]{babel}
\usepackage{graphicx}
\usepackage[export]{adjustbox}
\usepackage{amsmath}
\usepackage{amssymb}
\graphicspath{{/app/Quiz/Scripts/2023/ProgAll/15BinTree/}{/app/Quiz/Scripts/2021/Mod2021_3/BinTree/}{/app/Quiz/Scripts/2020/Mod2020_4/BinTree/}}
\setlength{\parindent}{0pt}
\setlength{\parskip}{6pt}
\pagestyle{empty}
\begin{document}
%begin
{\large\textbf{Контрольна робота}}

\textbf{Варіант \No\,10}\hfill \textit{ПІБ:}\,\underline{\hspace{60mm}}
\vspace{4mm}

%beginitem
1. Що буде виведено за виконання фрагменту коду?

%code
#include <iostream>

int g(int a, int b){
    int c = 45;
    if (a)
        c = 1;
    else if (b <= -1)
         return 4;
    else 
        c = 7;
    return c;
}

int main(){
    std::cout << g(3, 3);
    return 0;
}
%code
\vspace{5mm}
%enditem

%beginitem
2. 
try:
    res = 5.0j<="9.0"
    if isinstance(res, bool):
        print(f'bool;{res}')
    elif isinstance(res, int):
        print(f'int;{res}')
    elif isinstance(res, float):
        print(f'float;{res:.2f}')
    else:
        print('???')
except:
    print('Z')
\vspace{5mm}
%enditem

%beginitem
3. %goodpath:Scripts/2018/Mod2018_1/03-good.txt
Відмітити все, що є коректним оператором С++:
( 1 ) return x<=2;

( 2 ) double n.m;

( 3 ) cin>>x+3;

( 4 ) if (1>2) x=10;
\vspace{5mm}
%enditem

%beginitem
4. 

\begin{center}
	\adjustimage{max size={0.8\linewidth}{0.8\paperheight}}
	{../BinTree/trees_exp/35.png}
\end{center}
Указати послідовність міток вершин дерева, зображеного на рис., що утвориться в результаті обходу дерева в оберненому порядку. Мітки записувати через пробіл.\par Обхід дерева починається з кореня (на рис. це верхня вершина).
%answer:35 оберненому
\vspace{5mm}
%enditem

%beginitem
5. 
def f(n):
    print('f', end='')
    return n<=0

try:
    print(f(3) or f(-1))
except:
    print('ERR')
\vspace{5mm}
%enditem

%beginitem
6. 
b=['0','1','2','3','4']; b[-8:-8]='12'; print(b)\vspace{5mm}
%enditem

%beginitem
7. Що буде виведено за виконання фрагменту коду?

%code
print(7//7/7*7)
%code
\vspace{5mm}
%enditem

%beginitem
8. 
try:
    print(1, end=';')
    print(int('4'), end=';')
    print(7, end=';')
except KeyboardInterrupt:
    print(0, end=';')
except Exception:
    print(9, end=';')
except:
    print('Z', end=';')
else:
    print(5, end=';')
finally:
    print(6)
\vspace{5mm}
%enditem

%beginitem
9. %goodpath:Scripts/2019/Mod2019_1/Lex/03-good.txt
Відмітити все, що є коректним оператором С++:
( 1 ) return 'a'<'c'

( 2 ) print(sep=":","qwe")

( 3 ) x=*2

( 4 ) float(2)/87
\vspace{5mm}
%enditem

%beginitem
10. %goodpath:Scripts/2018/Mod2018_1/01-good.txt
Відмітити все, що є однією правильною лексемою:
( 1 ) a2f

( 2 ) 'qwerty'

( 3 ) **

( 4 ) true
\vspace{5mm}
%enditem

%end
\newpage
\end{document}

